% Ampliación original a partir de 46 diapositivas de fractura, ULE.
% Datos de ejercicios hipotéticos; no se reproducen figuras docentes.
\chapter{Fractura: ampliación basada en las diapositivas del profesor}
\label{ch:fractura-clase}
\objetivos{Entender clivaje, transición dúctil--frágil, Inglis y Griffith; aplicar los campos de punta y seleccionar factores geométricos para defectos, placas y probetas; comprobar unidades y validez antes de emitir conclusiones.}

\section{Mapa del capítulo docente}
Las diapositivas 4--15 presentan fractura, fallo histórico, mecanismos y transición térmica; 17--25 introducen cohesión, concentración y Griffith; 27--32 tratan los modos y la relación $K_I$--$G_c$; 33--36 recogen geometrías; 37--43 abordan tenacidades, diseño e inspección. Esta ampliación complementa el capítulo~\ref{ch:03} y el formulario; el cuaderno contiene ejercicios nuevos que reproducen los métodos, no los enunciados.

\section{Micromecanismos y temperatura}
En rotura dúctil metálica se identifican nucleación, crecimiento y coalescencia de microcavidades, con deformación plástica. El clivaje frágil transgranular puede seguir planos cristalográficos; la fractura intergranular puede estar vinculada a precipitación de fases frágiles, hidrógeno, corrosión, degradación ambiental o cavitación. La superficie de fractura ayuda a diagnosticar, pero no determina el mecanismo sin información de servicio.

Los casos docentes de buques Liberty y la lanzadera Challenger ilustran cómo interactúan material, temperatura, geometría y decisiones de operación. Son fenómenos históricos multicausales: no convertir una explicación simplificada de una diapositiva en diagnóstico exclusivo de un accidente. La transición dúctil--frágil se estudia mediante ensayos Charpy/Izod, energía absorbida y temperatura. Su localización depende de método, probeta, velocidad, orientación y criterio. Muchos aceros BCC son sensibles a la temperatura; no se afirma por ello que todo material FCC sea inmune a cualquier rotura frágil.

\section{Tensión de cohesión y defectos}
La separación de enlaces atómicos da una resistencia \emph{teórica} del orden de $E/10$ a $E/20$ para ciertos modelos, o una escala elemental $E/(2\pi)$. Las diapositivas comparan con resistencias experimentales que pueden ser mucho menores por defectos y mecanismos microestructurales. Estas escalas no son valores característicos certificados de una aleación.

Para una elipse central en placa infinita sometida a tracción perpendicular, si $a_e$ es \emph{semieje mayor}, $b_e$ semieje menor y $\rho=b_e^2/a_e$ el radio de punta, el modelo clásico de Inglis establece:
\begin{equation}
K_t=1+\frac{2a_e}{b_e}=1+2\sqrt{\frac{a_e}{\rho}}.
\label{eq:clase-inglis}
\end{equation}
La diapositiva 20 escribe en el radical una longitud $2a$ dibujada como longitud \emph{total}, mientras que esta definición clásica usa el semieje $a_e$. La discrepancia de convención afecta al factor $\sqrt2$: \emph{no identificar ambas expresiones automáticamente}. En un ejercicio con fórmula prescrita se deben explicitar las convenciones y, si son incompatibles, dejar constancia. Una grieta de punta aguda se estudia mejor con $K_I$ que con una tensión máxima elástica infinita.

\section{Derivación del criterio energético de Griffith}
Una grieta central pasante de longitud $2a$ en placa elástica ancha de espesor $B$ libera, en la idealización considerada, energía potencial
\[
\Delta\Pi=-\frac{\pi a^2 B\sigma^2}{E'},\qquad
E'=\begin{cases}E&\text{tensión plana},\\ E/(1-\nu^2)&\text{deformación plana}.\end{cases}
\]
Crear dos superficies de área $2aB$ cada una cuesta $4aB\gamma_s$, con energía superficial $\gamma_s$ en $\mathrm{J/m^2}$. La condición de equilibrio de la energía total da:
\begin{equation}
-2\frac{\pi aB\sigma_c^2}{E'}+4B\gamma_s=0
\quad\Longrightarrow\quad
\sigma_c=\sqrt{\frac{2E'\gamma_s}{\pi a}}.
\label{eq:clase-griffith}
\end{equation}
La expresión con $E$ de la diapositiva 25 corresponde a tensión plana. En materiales con disipación, la resistencia a propagación no se reduce necesariamente a dos energías de superficie libres. Para MFEL isotrópica:
\[
G=\frac{K_I^2}{E'},\qquad
G_c=\frac{K_{\mathrm{mat}}^2}{E'},\qquad
[G]=\mathrm{J/m^2},\quad[K]=\mathrm{MPa}\sqrt{\mathrm m}.
\]
La elección de $K_{\mathrm{mat}}$ depende de estado de esfuerzos, espesor y validez del ensayo.

\section{Campo del modo I y su alcance}
Con eje $x$ sobre la grieta, $y$ perpendicular y polar $(r,\theta)$ desde la punta, las tensiones dominantes en MFEL son:
\begin{align}
\sigma_{xx}&=\frac{K_I}{\sqrt{2\pi r}}\cos\frac{\theta}{2}
\left(1-\sin\frac{\theta}{2}\sin\frac{3\theta}{2}\right),\\
\sigma_{yy}&=\frac{K_I}{\sqrt{2\pi r}}\cos\frac{\theta}{2}
\left(1+\sin\frac{\theta}{2}\sin\frac{3\theta}{2}\right),\\
\tau_{xy}&=\frac{K_I}{\sqrt{2\pi r}}\sin\frac{\theta}{2}
\cos\frac{\theta}{2}\cos\frac{3\theta}{2}.
\end{align}
Sobre el eje frontal $\theta=0$, $\sigma_{xx}=\sigma_{yy}=K_I/\sqrt{2\pi r}$ y $\tau_{xy}=0$. Los signos y nombres de componentes dependen de cómo se dibujen ejes y ángulos; la diapositiva 29 no debe mezclarse con otra convención sin comprobarlos. La singularidad solo describe un término asintótico fuera de la zona plástica y a distancia pequeña frente a dimensiones de la pieza; no predice tensiones infinitas reales.

\section{Deformación plana y criterio de validez}
En tensión plana $\sigma_z=0$. En deformación plana isótropa lineal, $\varepsilon_z=0$ y $\sigma_z=\nu(\sigma_x+\sigma_y)$. Un cribado de dimensiones para una tenacidad de deformación plana es:
\begin{equation}
B,\ a,\ W-a\gtrsim2.5\left(\frac{K_{IC}}{\sigma_y}\right)^2.
\label{eq:clase-validez}
\end{equation}
La diapositiva 31 presenta solo la condición sobre $B$. Esta no demuestra por sí sola una medición válida de $K_{IC}$: faltan longitud de grieta, ligamento, cargas, geometría de probeta y otros criterios normalizados. No se equiparan indiscriminadamente $K_c$ aparente de una chapa delgada y $K_{IC}$ válido de probeta gruesa.

\section{Factores de intensidad de las geometrías del profesor}
Las diapositivas 33--36 emplean $C$ como nombre alternativo de factor geométrico $Y$ \emph{en las expresiones $K_I=C\sigma\sqrt{\pi a}$}. Se diferencia de las funciones adimensionales $f$ de probetas CT o SENB, que normalizan la \emph{carga} según su propia geometría.

Para grieta \emph{central} con longitud total $2a$ en placa de ancho completo $W$, una aproximación de anchura finita es:
\begin{equation}
Y_{\mathrm c}=\sqrt{\sec\left(\frac{\pi a}{W}\right)}
\quad(0<2a/W<1).
\label{eq:clase-central}
\end{equation}
La diapositiva 33 incluye asimismo un ajuste polinómico y una aproximación denominada Dixon. Son expresiones alternativas y no se deben intercambiar fuera del dominio de calibración. Al tender $a/W\to0$, $Y_{\mathrm c}\to1$. La tensión nominal debe definirse igual que en la solución.

Para grieta \emph{simple de borde} de longitud $a$ y razón $x=a/W$, el ajuste escrito en la diapositiva 33 es:
\begin{equation}
Y_{\mathrm b}(x)=1.12-0.231x+10.55x^2-21.72x^3+30.39x^4,
\qquad 0\le x\lesssim0.6.
\label{eq:clase-borde}
\end{equation}
El valor $1.12$ es el límite de grieta corta, no una corrección fija para grietas largas. Las dos grietas de borde opuestas requieren el factor distinto de la diapositiva 34, que depende del ligamento remanente.

Una \emph{grieta circular interna} idealizada (circular en su plano, no agujero redondo) tiene $K_I=(2/\pi)\sigma\sqrt{\pi a}$ con $a$ radio. La grieta elíptica interna depende del punto del frente y de la proporción entre ejes; no asignar $Y=1$ sin justificarlo. Tampoco confundir $K_t$ del agujero redondo de la diapositiva 21 con intensidad de grieta.

Una probeta compacta CT utiliza una normalización específica:
\[
K_I=\frac{P}{B\sqrt W}f(a/W),\quad
f(x)=\frac{(2+x)(0.886+4.64x-13.32x^2+14.72x^3-5.60x^4)}
{(1-x)^{3/2}}.
\]
La fórmula se aplica a su geometría y rango de calibración; no basta introducir en ella la carga de una chapa. La probeta de flexión SENB exige el vano y el factor de la figura de la diapositiva 35; el doble voladizo DCB se relaciona con rigidez de brazos, carga, grieta y estado plano (diapositiva 36). Para ambos se explican y se fijan primero los significados de $P,B,W,a,S,h$ del dibujo; no se copia el coeficiente de CT por analogía.

\section{Magnitudes de diseño, detectar y mantener}
Para $Y$ \emph{constante}, en una grieta central o de borde con $a$ definido coherentemente:
\begin{equation}
\sigma_{\mathrm{adm}}=\frac{K_{\mathrm{mat}}}{Y\sqrt{\pi a}},
\qquad a_c=\frac1\pi\left(\frac{K_{\mathrm{mat}}}{Y\sigma_{\max}}\right)^2.
\label{eq:clase-disenyo}
\end{equation}
Si $Y=Y(a/W)$, el tamaño crítico exige resolver la ecuación implícita; la longitud total de grieta central será $2a_c$. La diapositiva 40 contiene un rótulo dimensional inconsistente: $K_{IC}$ tiene unidad $\mathrm{MPa}\sqrt{\mathrm m}$, no $\mathrm{MPa\,m^{-2}}$; $G_c$ y $\gamma_s$ se expresan en $\mathrm{J/m^2}$; $Y$ y $K_t$ son adimensionales.

Un plan tolerante al daño considera solicitación adversa, tamaño crítico, crecimiento por fatiga y capacidad real de END; un equipo nominalmente capaz de ver un tamaño no demuestra detección del cien por cien. Las inspecciones son periódicas por riesgo y datos de crecimiento, y el cálculo académico no es instrucción de aeronavegabilidad.

\repaso{¿Por qué un $K_t$ conocido no sustituye a $Y$?}{$K_t$ compara máxima tensión elástica y tensión nominal de una entalla; $Y$ corrige la intensidad del campo singular de una grieta. La igualdad de dimensiones no implica igualdad de modelos.}
\repaso{Una grieta central mide $8\,\mathrm{mm}$ de extremo a extremo. ¿Qué se utiliza en $K_I=Y\sigma\sqrt{\pi a}$?}{Se usa $a=4\,\mathrm{mm}=0.004\,\mathrm m$ si el factor corresponde a longitud total $2a$; otra geometría exige redefinir el parámetro.}
\resumen{El procedimiento completo es: geometría y convención de longitudes, solicitación y $Y$, cálculo de $K_I$, comparación con resistencia aplicable, validez de MFEL, unidades, interpretación y límite de la conclusión.}
\fuentes{Pablo Rodríguez González, \textit{Mecánica de la fractura}, presentación docente ULE de 46 diapositivas facilitada para esta revisión; páginas 4--15, 19--25, 27--36, 37--43. Fundamentos contrastados: Roylance, \textit{Introduction to Fracture Mechanics} \cite{MIT-FRAC}. Se distingue el texto visible de la presentación de las correcciones de convención y advertencias inferidas.}
